\section{Основные правила игры}

\subsection{Начало игры}

\subsubsection{Игровой набор}

Для игры используются 136 специальных костей, которые чаще называют \textit{тайлами}. Есть 34 вида тайлов, и каждый из них встречается в наборе ровно четыре раза (34 х 4 = 136). На одной стороне тайла находится значимое изображение, другая же, рубашка, у всех тайлов одинакова.

\textbf{Тайлы мастей}:

\noindent\begin{tabular}{ C{3.5cm}ccccccccc }
	\toprule
	Масть/номинал & 1 & 2 & 3 & 4 & 5 & 6 & 7 & 8 & 9 \\
	\midrule
	Ман \newline & \mahjong{1m} & \mahjong{2m} & \mahjong{3m} & \mahjong{4m} & \mahjong{5m} & \mahjong{6m} & \mahjong{7m} & \mahjong{8m} & \mahjong{9m} \rule[1ex]{0pt}{7ex} \\
	\midrule
	Пин \newline & \mahjong{1p} & \mahjong{2p} & \mahjong{3p} & \mahjong{4p} & \mahjong{5p} & \mahjong{6p} & \mahjong{7p} & \mahjong{8p} & \mahjong{9p} \rule[1ex]{0pt}{7ex} \\
	\midrule
	Соу \newline & \mahjong{1s} & \mahjong{2s} & \mahjong{3s} & \mahjong{4s} & \mahjong{5s} & \mahjong{6s} & \mahjong{7s} & \mahjong{8s} & \mahjong{9s} \rule[1ex]{0pt}{7ex} \\
	\bottomrule
\end{tabular}

Тайлы мастей с номиналами 1 и 9 называются \textit{терминальными}, остальные тайлы мастей называются \textit{простыми} или \textit{срединными}.
\vspace{1cm}

\textbf{Благородные тайлы} (обратите внимание, порядок важен):

\noindent\begin{tabular}{ cC{2cm}C{2cm}C{2cm}C{2cm} }
	\toprule
	 \rule[0ex]{0pt}{7ex} Ветра & \mahjong{1z} \newline Восток & \mahjong{2z} \newline Юг & \mahjong{3z} \newline Запад & {\mahjong{4z} \newline Север} \\
\end{tabular}
\noindent\begin{tabular}{ cC{2cm}C{2cm}C{2cm} }
	\toprule
	 \rule[0ex]{0pt}{7ex} Драконы & \mahjong{5z} \newline Белый & \mahjong{6z} \newline Зеленый & \mahjong{7z} \newline Красный \\
	\bottomrule
\end{tabular}

\vspace{1cm}
Для подсчета очков в игре используются счетные палочки следующих номиналов:

\noindent\begin{tabular}{cc}
	\toprule
	\includegraphics{tenbo10000} & 10000 очков \\
	\includegraphics{tenbo5000} & 5000 очков \\
	\includegraphics{tenbo1000} & 1000 очков \\
	\includegraphics{tenbo100} & 100 очков \\
	\bottomrule
\end{tabular}

Иногда можно встретить в наборах палочки номиналом 500 очков, обычно они выглядят так же, как палочки на 100 очков, но помеченные особым образом (например, окрашенные в зеленый цвет).

Общее количество очков в начале игры у каждого игрока равно 30000.
Также в игре используются:

\begin{itemize}
	\item Индикатор первого дилера --- небольшая пластинка с изображением восточного ветра (\textnihon{東}) с одной стороны и южного ветра (\textnihon{南}) с другой;
	\item Два шестигранных кубика с цифрами от 1 до 6
\end{itemize}

%TODO добавить фото индикатора

Помимо упомянутого, можно также использовать специальный компас с указанием ветров, который кладется в центр стола на время раздачи и убирается на время перемешивания тайлов.

\begin{figure}[H]
	\centering
	\includegraphics[width=8cm]{compass.jpg}
	\caption{Специальный компас с указанием ветров}
\end{figure}

\subsubsection{Подготовка к игре}

В риичи-маджонг играют вчетвером. Для игры используют небольшой квадратный стол (\sim75х75 см), с каждой стороны которого садится по игроку. Каждому месту за столом присваивается условная сторона света, и расположены они в порядке тайлов ветров против часовой стрелки: восток (\textnihon{東}), юг (\textnihon{南}), запад (\textnihon{西}), север (\textnihon{北}), т. е. порядок "неправильный" и отличается от настоящего расположения сторон на карте мира\footnote{Порядок ветров соответствует расположению сторон света на карте звездного неба.}. Выбор мест производится по договоренности или по жребию. Ветер, соответствующий месту игрока, называется \textit{ветром места}.

Игрок, сидящий на востоке, становится первым \textit{дилером} в игре и рядом с ним в начале игры кладется индикатор раунда. Обычно дилер выбирает место за столом, остальные рассаживаются относительно него согласно вытянутым ветрам.

Игра делится на раунды --- восточный (первый) и южный (второй). Раунды делятся на раздачи.

Ветер, соответствующий текущему раунду, называется \textit{ветром раунда}. Раунд по умолчанию состоит из четырех раздач (хотя в отдельных случаях, оговоренных правилами, могут назначаться дополнительные раздачи), и каждую новую раздачу происходит сдвиг сторон света за столом на одно место против часовой стрелки: игрок, в первой раздаче бывший на юге, во второй оказывается дилером, в третьей --- севером и т. д, сами игроки при этом не пересаживаются. Смена раунда происходит, когда дилерство возвращается к игроку, бывшему дилером в первой раздаче. Индикатор первого дилера все время остается лежать рядом с дилером первой раздачи; во время восточного раунда он повернут вверх стороной, где нарисован иероглиф востока, а с началом южного переворачивается. Для указания на дилера текущей раздачи с его стороны стола после подготовки к раздаче кладутся кубики. В случае, если используется специальный компас, кубики можно размещать в специально предусмотренном для них месте компаса на время разыгрывания раздачи.

Перед началом каждой раздачи игроки кладут все тайлы на стол рубашкой вверх и тщательно их перемешивают. После этого каждый игрок строит перед собой "стену" высотой в 2, шириной в 1 и длиной в 17 тайлов рубашкой вверх, размещая ее так, чтобы в итоге четыре участка стены от разных игроков образовали замкнутый квадрат.

%TODO фото собранной стены

После построения стены дилер однократно бросает кубики, озвучивает выпавшее на кубиках число и убирает кубики в правый нижний угол стола. Бросать кубики следует с высоты 10-15 см от поверхности стола, чтобы гарантировать честное случайное число, выпадающее на кубиках.

Затем дилер отсчитывает против часовой стрелки столько сторон квадрата, сколько выпало на кубиках суммарно, начиная со своей стороны (на рисунке показан пример для выпавшей суммы 8).

\begin{figure}[H]
	\centering
	\includegraphics[width=15cm]{table-deal-start.png}
	\caption{Определение разлома стены}
\end{figure}

Затем игрок, на чью сторону указал дилер (в примере выше --- север), отсчитывает от правого конца своей стороны стены столько тайлов, какое число выпало на кубиках до этого (в примере --- 8) и немного отодвигает их от соседних тайлов вправо, создавая разлом стены. После этого он отсчитывает от разлома в обратную сторону 7 тайлов и немного отодвигает получившуюся группу из 14 тайлов (см. следующий рисунок). Эта группа называется \textit{мертвой стеной} и не разбирается\footnote{За исключением одного особого случая при объявлении кана, см. далее} в процессе игры. Оставшаяся часть стены называется \textit{живой стеной} и используется во время раздачи для получения игроками тайлов. В ходе игры стена разбирается по одному тайлу (сначала берется верхний тайл в стопке, потом нижний) по часовой стрелке, заканчивая нижним тайлом последней стопки перед мертвой стеной. Стена считается непрерывной, и переход через углы не имеет никакого значения как при формировании мертвой стены, так и при последующем разборе.

Третий верхний тайл от конца мертвой стены переворачивается рубашкой вниз (на рисунке это 6 ман) и служит в данной раздаче индикатором доры --- тайла, дающего дополнительные очки при его наличии в руке. Подробно о дорах и принципе работы индикаторов рассказано в разделах далее.

\begin{figure}[H]
	\centering
	\includegraphics[width=16cm]{table-deal-start-2.png}
	\caption{Мертвая и живая стены}
\end{figure}

После открытия индикатора доры игроки начинают набор своих стартовых рук со стены\footnote{Допускается открыть индикатор доры после того как руки набраны, но нужно это сделать до первого сброса дилера}. Начиная с дилера по очереди против часовой стрелки игроки начинают брать от начала живой стены по четыре тайла (две стопки по 2), пока у каждого не окажется 12 тайлов.
После этого дилер берет один крайний тайл со стены и тайл через один от него, далее игрок на юге берет нижний тайл, игрок за западе - верхний тайл перед тем, который уже взял дилер, и игрок на севере забирает один оставшийся тайл.

Таким образом, в начале раздачи дилер имеет 14 тайлов, а остальные --- по 13, а стена начинается с одной неполной стопки, имеющей только один нижний тайл. Тайлы своих рук игроки расставляют перед собой в ряд вертикально изображением к себе, так, чтобы остальные видели только их рубашки. Сортировка тайлов в руке оставляется на усмотрение игрока, однако в случае показательных игр (например, при игре на камеру) рекомендуется сортировать тайлы в руке.

\begin{figure}[H]
	\centering
	\includegraphics[width=16cm]{table-deal-start-3.png}
	\caption{Порядок разбора стены, фото}
\end{figure}

Обратите внимание на подписи тайлов --- они показывают, какой игрок должен их взять (верхние тайлы стены подписаны прямо на рубашке, нижние --- рядом). Предположим, что сейчас идет первая раздача восточного раунда: на индикаторе первого дилера указан восток, а кубики, как индикатор текущего дилера, лежат рядом с ним, т. е. первый дилер является дилером в этой раздаче.

\subsection{Ход раздачи}

Раздача начинается с хода дилера, который сбрасывает с руки один из своих 14 тайлов, кладя его рубашкой вниз в области сброса недалеко от центра стола; затем ход переходит к следующему игроку против часовой стрелки, который берет первый тайл из оставшейся части живой стены и также сбрасывает один тайл из своей руки (можно сбрасывать и тайл, только что взятый со стены). Затем ход переходит к следующему игроку. Так раздача продолжается до тех пор, пока кто-либо не соберет руку и объявит победу, или же пока в живой стене не закончатся тайлы. Ходы четырех игроков, начиная от дилера, образуют круг раздачи. После окончания круга ход вновь переходит к дилеру, и начинается следующий круг. Группы сброшенных тайлов в центре стола называются сбросами или \textit{дискардами}. Каждый игрок имеет свой отдельный дискард. Тайлы следует выкладывать в сброс слева направо тремя рядами по 6 штук (начиная с ближнего к центру стола и далее по направлению к себе); четвертый же ряд обычно не начинают, а вместо этого продолжают третий. Цель каждого игрока в раздаче --- собрать в руке четыре сета и одну пару.При объявлении кем-либо победы с готовой рукой раздача заканчивается. Цель всей игры --- набрать как можно больше очков.

Сет --- это определенная комбинация из трех или четырех тайлов. Они бывают трех видов:

\begin{itemize}
	\item Последовательность из трех тайлов одной масти подряд (по типу 1-2-3, 3-4-5 и т. п.) --- \textit{чи}. Последовательности можно собирать только из тайлов мастей, и в них не должно быть перехода через девятку: 8-9-1 и 9-1-2 не засчитываются как чи.
	\item Три одинаковых тайла --- \textit{пон};
	\item Четыре одинаковых тайла --- \textit{кан} (должен быть \textit{объявлен}, про объявления см. далее).
\end{itemize}

Один тайл не может одновременно входить в два сета или в сет и пару.

\begin{example}

Пример: в руке уже есть готовый пон из 4 ман, чи 7-8-9 ман и пара красных драконов, но последовательность 3-4-5-6-7 пин не засчитывается за два чи: для завершения этой формы в два сета требуется получить еще один тайл (2, 5 либо 8 пин).
\vspace{5mm}

\mahjong{444789m 34567p 77z}

\end{example}

\subsection{Принцип яку}

Собрать на руке четыре сета и пару --- это необходимое условие для победы, но \textit{не достаточное}. Для того, чтобы иметь право объявить победу, в руке должны быть соблюдены какие-либо условия, дающие \textit{стоимость} руке. В риичи-маджонге такие условия называются \textit{яку} (хотя иногда их называют \textit{комбинациями}, что не полностью верно, поскольку помимо определенных сочетаний тайлов яку могут давать и некоторые игровые ситуации). Яку могут суммироваться (если в описании яку не указано обратное), таким образом давая еще большую стоимость руки. Стоимость яку измеряется в \textit{ханах}. Полный список яку с примерами и описаниями приведен в конце документа.

\begin{example}

	Пример: в руке ниже есть три красных дракона, и это дает яку, которое называется \textit{якухай}, и оно стоит 1 хан. С такой рукой можно объявить победу.
	\vspace{5mm}

	\mahjong{444789m 4567p 777z}

\end{example}

\subsection{Объявления в игре}

Помимо взятий со стены и сбросов во время раздачи игроки могут делать различные \textit{объявления}. Основных видов объявлений два: это объявления на тайлы (объявления сетов) и объявления победы. Помимо этого, существует особое объявление риичи, о нем рассказывается в соответствующем разделе.

 Объявить сет с чужим сброшенным тайлом можно только в момент его сброса другим игроком и до взятия со стены тайла следующим игроком --- после этого тайл в дискарде считается вышедшим из игры, и больше его использовать для объявлений нельзя. Также нельзя делать никакие объявления, кроме объявления победы, на последний сбрасываемый тайл за раздачу, когда в живой стене не остается тайлов.

Рука, в которой есть объявленные сеты с чужими тайлами, называется \textbf{открытой}, а рука без них --- \textbf{закрытой}. Сет, образованный любым объявлением, считается фиксированным, и изменять его нельзя. Сбросы разрешается делать только из оставшейся закрытой части руки.

\subsubsection{Объявление "чи"}

Чи --- это сет, образованный любыми тремя тайлами одной и той же масти в непрерывающейся последовательности (например, 4-5-6). Если в руке есть два любых тайла, которые могут входить в одну последовательность, при сбросе противником третьего тайла этой последовательности объявлением "чи" можно дополнить эти два тайла до чи. Чи можно объявлять только на тайлы, сброшенные предыдущим игроком по ходу игры (т.е. сидящим слева). Последовательность действий при объявлении "чи" :

\begin{itemize}
	\item После сброса игроком слева тайла в дискард нужно сказать "чи";
	\item Затем следует выложить два тайла последовательности из руки плашмя рубашкой вниз справа;
	\item После этого нужно взять тайл из дискарда противника и дополнить им последовательность, положив взятый тайл боком слева от двух открытых тайлов;
	\item Наконец, следует осуществить сброс тайла из руки.
\end{itemize}

Два последних действия можно выполнять в любом порядке.

Взятый при "чи" тайл всегда кладется слева от двух других, даже если при этом нарушается порядок последовательности. Если с рукой из примера выше, где уже взят пон на 2 ман, взять на следующем круге чи при сбросе игроком слева 4 пин, она будет выглядеть так:

\mahjong{345m 7m 4888s} \hfill \mahjong{4'23p-22'2m}

Каждый следующий объявленный сет кладется слева либо сверху от предыдущего.

\subsubsection{Объявление "пон" }

\textbf{Пон} --- это три одинаковых тайла. Объявление пона позволяет дополнить имеющиеся в руке два одинаковых тайла третьим таким же тайлом, когда его сбрасывает другой игрок. Пон можно объявить на тайл, сброшенный любым из трех противников. Последовательность действий при объявлении "пон" аналогична последовательности при объявлении "чи":

\begin{itemize}
	\item После сброса нужного тайла противником в дискард нужно сказать "пон";
	\item Затем следует выложить свою пару тайлов плашмя рубашкой вниз справа от руки;
	\item После этого нужно взять тайл из дискарда противника и положить его боком слева, между или справа от тайлов пары в зависимости от того, с игрока слева, напротив или справа был взят пон.
	\item Наконец, после объявления сбрасывается тайл из закрытой части руки, и ход переходит к следующему игроку от взявшего пон против часовой стрелки, т.е. объявление нарушает порядок ходов, прерывая круг.
\end{itemize}

Два последних действия можно выполнять в любом порядке.

\newpage

Выглядит это так (поны с игроков слева, напротив и справа соответственно):

\mahjong{4'44m-55'5m-666'm}

\begin{example}

Например, рука

\vspace{5mm}

\mahjong{22345m 237p 4888s 2z}

\vspace{5mm}

после взятия 2 ман с игрока напротив будет выглядеть так:

\vspace{5mm}

\mahjong{345m 237p 4888s 2z} \hfill \mahjong{22'2m}

\end{example}

\subsubsection{Объявление "кан"}

"Кан" --- объявление, позволяющее создавать в руке каны, одноименные сеты из четырех одинаковых тайлов. Если в закрытой части руки есть четыре одинаковых тайла, но кан не был объявлен, они не интерпретируются как кан (поскольку иначе в руке не хватит тайлов для четырех сетов и пары), но могут интерпретироваться любым другим способом. Например, в руке ниже кан не объявлен, и поэтому четыре тайла 5 ман не считаются за кан, но могут считаться за пон и часть чи 5-6-7, либо за пару, часть чи и свободный тайл и т.д.

\mahjong{505567m 12p 99s 3356z}

Кан можно объявлять тремя различными способами.

\textbf{Закрытый кан (\textit{анкан})}

Если игрок имеет в закрытой части руки четыре одинаковых тайла, он в любой свой ход после взятия тайла со стены, но до сброса, может сказать "кан" и выложить эти четыре тайла плашмя справа, положив два крайних тайла рубашкой вверх, а два средних --- рубашкой вниз (крайние тайлы следует предварительно показать противникам). С закрытым каном 5 ман из руки выше это будет выглядеть так:

\mahjong{67m 12p 99s 3356z} \hfill \mahjong{X50mX}

Рука с закрытым каном продолжает считаться закрытой --- это единственный объявляемый сет-исключение.

\textbf{Открытый кан (\textit{дайминкан})}

Если у игрока в закрытой части руки есть три одинаковых тайла, и любой противник сбрасывает четвертый, он может объявить на него "кан", действуя точно так же, как и при объявлении "пон". Тайлы кана кладутся плашмя справа таким же образом, как и после "пона". Если кан был взят с игрока напротив, можно повернуть боком любой из двух средних тайлов. Пример: если игрок из примеров выше, до этого взявший пон и чи, берет на следующем после чи круге кан 8 соу с игрока справа, его рука будет выглядеть как на схеме ниже.

\mahjong{345m 7m} \hfill \mahjong{8888's-4'23p-22'2m}

\textbf{Дополненный кан (\textit{шоминкан})}

Если у игрока есть открытый пон и четвертый такой же тайл в закрытой части руки, в любой свой ход после взятия со стены и до своего сброса он может дополнить пон до кана объявлением "кан". Четвертый тайл кладется боком сверху от тайла, лежавшего боком в поне. Такой кан, как и закрытый, нельзя объявлять сразу после своего объявления "чи" или "пон" (но можно после другого любого кана). Дополненный кан считается открытым.

\mahjong{345m 7m} \hfill \mahjong{888'-4'23p-22"2m}

Закрытый и дополненный каны нельзя объявлять сразу после своего объявления чи или пона (но можно после другого кана).

После объявления любого кана в закрытой части руки начинает не хватать тайлов для ее завершения, и поэтому игрок сразу же должен взять с мертвой стены тайл замены и только потом совершить сброс. Это единственная ситуация в игре, когда берутся тайлы с мертвой стены.

Кроме того, при объявлении каждого кана открывается один индикатор \textit{кандоры}: переворачивается соседний с индикатором доры верхний тайл с "длинной" стороны мертвой стены. Новый индикатор открывается \textit{сразу} после объявления и до взятия игроком тайла замены с мертвой стены.

При дополнении пона до кана, дополняющий тайл может быть затребован для победы, в таком случае новый индикатор доры не открывается, так как кан фактически не был объявлен.

В мертвой стене всегда должно сохраняться 14 тайлов. Поэтому при взятии тайла замены последний тайл живой стены становится частью мертвой стены.

Порядок действий при объявлении кана:

\begin{itemize}
	\item Четко объявить вслух "кан"
	\item Открыть тайлы из руки (четыре для закрытого кана, три для открытого кана и один для дополненного кана) и переместить их вправо от руки;
	\item Перевернуть новый индикатор кандоры;
	\item Взять тайл замены из мертвой стены;
	\item Сбросить тайл из руки и продолжить ход.
\end{itemize}

Победу на дополняющем тайле (\textit{чанкан}) следует объявлять не ранее, чем четвертый тайл будет доложен в кан и не позднее, чем игрок, объявивший кан, осуществит сброс тайла в дискард либо объявит еще один кан.

Всего за раздачу может быть объявлено не более чем 4 кана: именно столько тайлов замены возможно взять с короткого конца мертвой стены, и столько индикаторов кандор можно открыть с длинного.

\subsubsection{Куикаэ}

Правило \textit{куикаэ} накладывает дополнительные ограничения на объявления сетов --- запрещает сброс такого же тайла после взятия сета, а также сброс тайла, ранее дополнявшего до последовательности объявляемые два тайла ("сдвиг чи").

%возможно надо выкинуть, оставив только в разделе про нарушения
% В случае нарушения данного правила никаких изменений в дискарде делать не следует (т.е. нельзя забрать тайл обратно в руку и скинуть другой тайл), при этом к игроку-нарушителю применяются санкции (подробнее о нарушениях и санкциях см. в разделе про судейство). Объявление сета, завершающееся сбросом некорректного тайла с точки зрения правила куикаэ, считается корректным.

\begin{example}

Пример:
\vspace{5mm}

\mahjong{2255m 5678s 333z} \hfill \mahjong{5'34m}
\vspace{5mm}

Сброс как 2 ман, так и 5 ман в данном случае недопустим из-за правила куикаэ. Аналогичным образом, после взятия пона не допускается сброс такого же тайла, какой был взят в пон.

\end{example}

\subsubsection{Объявления победы}

Когда игроку не хватает в руке до победы одного тайла, говорят, что у игрока \textit{темпай}\footnote{Также встречаются варианты "Игрок темпай", или "Игрок в темпае", все они означают одно и то же --- состояние руки, которой до победы не хватает одного тайла}. Если на руке темпай, игрок может выиграть раздачу, получив выигрышный тайл двумя способами:

\begin{itemize}
	\item Со стены в свой ход. Такая победа называется \textit{цумо}, и объявляется словом "цумо" после взятия выигрышного тайла.
	\item При сбросе другим игроком на его ходу. Такая победа называется \textit{рон} и объявляется словом "рон".
\end{itemize}

Если все выигрышные тайлы находятся в руке, то такой игрок не считается темпай.

\begin{example}

Пример темпая:
\vspace{5mm}

\mahjong{23455m 34p 456s} \hfill \mahjong{7'68m}
\vspace{5mm}

Выигрышные тайлы, завершающие руку в четыре сета и пару --- 2 и 5 пин. Игрок может объявить цумо, если один из них зайдет со стены, или рон, если какой-то из них сбросит любой из оппонентов.

\end{example}

После объявления цумо выигрышный тайл кладется рядом с рукой, а рука показывается игрокам.

После объявления рон выигрышный тайл остается в дискарде, а рука также показывается игрокам.

Объявление победы завершает раздачу. После показа руки производится подсчет очков и выплата за руку, и затем, если раздача не была последней, начинается следующая. Объявления победы, как и объявления сетов, не обязательны:

\begin{itemize}
	\item рон можно пропустить, или сделать на этот тайл другое объявление --- пон, кан или чи;
	\item цумо тоже можно не объявлять, сбросив любой тайл из руки после получения выигрышного и продолжив раздачу.
\end{itemize}

Несколько игроков в темпае могут объявить рон на один и тот же сброшенный тайл одновременно. В этом случае выплаты производятся каждому из победивших игроков.

\subsubsection{Приоритет объявлений}

Когда игроки одновременно делают различные объявления на один тайл, забирает его тот игрок, чье объявление приоритетнее:

\begin{itemize}
	\item Пон или кан приоритетнее чи;
	\item Рон приоритетнее, чем пон или кан (тайл при этом остается в дискарде).
\end{itemize}

Последний момент, когда можно объявить рон --- непосредственно до того, как следующий по ходу игры игрок осуществит сброс.

\subsubsection{Риичи}

Риичи --- особое объявление, дающее одноименное яку \textit{риичи} и таким образом позволяющее выиграть на любой закрытой руке, находящейся в темпае, даже если других яку в ней нет. Также при помощи риичи можно увеличить стоимость руки, если в ней есть другие яку.

Игрок может объявить риичи в любой свой ход во время сброса тайла в дискард, если соблюдаются все следующие условия:

\begin{itemize}
	\item На руке есть темпай и этот темпай остается после сброса последнего лишнего тайла;
	\item Рука полностью закрытая;
	\item До исчерпания живой стены осталось не менее четырех тайлов
\end{itemize}

Порядок объявления риичи следующий:

\begin{itemize}
	\item Игрок говорит "риичи";
	\item Затем он сбрасывает тайл в дискард и кладет его боком;
	\item После этого он кладет в центр стола над своим дискардом палочку в 1000 очков --- ставку риичи.
\end{itemize}

В случае, если другой игрок объявляет рон с тайла, на котором объявляется риичи, ставка в 1000 очков не считается сделанной и не забирается выигравшим игроком.

Рассмотрим дискард после объявления риичи:

\hspace{1.3cm} \includegraphics{tenbo1000}

\mahjong{4z1m7p6z7z2s}

\mahjong{2m4m3'p7s}

Над дискардом --- палочка в 1000 очков, в самом дискарде один тайл, сброшенный при объявлении, повернут боком. Если после объявления кто-то забирает сброшенный тайл для сета, риичи-игрок обязан повернуть боком тайл, сбрасываемый на следующем ходу, чтобы один повернутый боком тайл обязательно был в дискарде. После него все сбрасываемые тайлы выкладываются обычным образом, как 7 соу на рисунке.

Объявив риичи, игрок не имеет права менять состав своей руки, и должен сбрасывать все приходящие со стены тайлы до получения выигрышного. Возможна победа как со стены (цумо), так и с чужого дискарда (рон).

После объявления риичи остается возможность объявить закрытый кан, используя закрытый пон в руке и только что пришедший тайл со стены, но только если кан никак не изменяет ожидание и возможности интерпретации сетов в руке. Никакие другие сеты после риичи объявлять нельзя.

\begin{example}

	\mahjong{2399m333444555s}
	\vspace{5mm}

	В руке выше c объявленным риичи ни один из возможных канов 3, 4 или 5 со при заходе соответствующих тайлов объявлять нельзя, хоть они и не изменят ожидания. При объявлении такого кана исчезла бы одна из возможных интерпретаций сетов в руке: форма 333444555 может быть рассмотрена не только как три пона, но и как три чи 3-4-5, но при объявлении кана такая возможность интерпретации исчезает.

\end{example}

Ставка риичи --- это сумма в 1000 очков, вычитающаяся из счета игрока. Если игрок выигрывает после объявления риичи, она возвращается ему обратно вместе с очками, полученными за руку. Если же раздачу выигрывает другой игрок, он забирает эту тысячу очков себе. В случае ничьей палочка остается лежать на столе до первой победы.

В случае двойного или тройного рона, все риичи-ставки, которые выиграли, возвращаются объявившим их игрокам. Лежащие на столе и невыигравшие ставки отдаются игроку, выигравшему первому по ходу игры от проигравшего. Если одна из победных рук не может быть засчитана по каким-то причинам (например, она является некорректной, т.к. в ней нет четырех сетов и пары), засчитываются все другие победы по рон, которые оказались корректными.

\subsection{Фуритен и временный фуритен}

Правило \textit{фуритен} , или правило упущенного сброса, запрещает игроку объявлять рон, если хотя бы один из тайлов, завершающих руку до выигрышной формы, лежит у него в дискарде. Учитываются также тайлы, затребованные из дискарда для сетов другими игроками.

\begin{example}

	Например, рассмотрим игрока со следующей рукой:
	\vspace{5mm}

	\mahjong{33456678m 23p 678s}
	\vspace{5mm}

	Предположим, дискард у игрока следующий:
	\vspace{5mm}

	\mahjong{1p 4z 7z 2s 4s 9m}
	\vspace{5mm}

	В данном случае игрок не может объявить рон ни при сбросе противником 1 пин, ни при сбросе 4 пин, потому что в его дискарде есть "упущенный" выигрышный тайл 1 пин. При этом ограничений на объявление цумо нет. Если игрок перестроит свою руку так, чтобы сброшенные в дискард тайлы не завершали выигрышную форму, фуритен перестанет действовать, и он снова сможет объявлять рон.

\end{example}

\textit{Временный фуритен} --- правило, которое не позволяет игроку объявлять рон, если на текущем круге считая от хода игрока уже был сброшен какой-либо из тайлов, завершающих его руку до выигрышной формы. Действие временного фуритена заканчивается после следующего сноса игрока. Исключение --- игрок, объявивший риичи, остается во временном фуритене до конца раздачи.

\subsection{Дора}

Дора --- это тайл, наличие которого в руке увеличивает ее стоимость на 1 хан. Однако дора --- это своего рода бонус, который сам по себе не дает яку, поэтому с рукой без яку нельзя объявить победу даже если в ней есть доры. Дора в каждой раздаче определяется индикатором, который открывается в мертвой стене перед началом раздачи. Дорой считается следующий тайл после индикатора по порядку:

\begin{itemize}
	\item Тайлы мастей: 1\rightarrow
	2\rightarrow
	3\rightarrow
	4\rightarrow
	5\rightarrow
	6\rightarrow
	7\rightarrow
	8\rightarrow
	9\rightarrow
	1 (дора только в той же масти, что и индикатор);
	\item Ветра: \textnihon{東} (восток)\rightarrow
	\textnihon{南} (юг)\rightarrow
	\textnihon{西} (запад)\rightarrow
	\textnihon{北} (север)\rightarrow
	\textnihon{東} (восток);
	\item Драконы: красный\rightarrow
	белый\rightarrow
	зеленый\rightarrow
	красный.
\end{itemize}

Помимо "обычных" дор существуют дополнительные. За каждый экземпляр любой из них дается дополнительный хан, как и за обычную дору:

\begin{itemize}
	\item \textit{Кандоры}. При объявлении любого кана в мертвой стене открывается один индикатор кандоры. Сколько в раздаче объявлено канов, столько и открыто индикаторов кандор. Индикаторы кандоры работают точно так же, как и обычные.
	\item \textit{Урадоры}. При победе после объявления риичи, тайл, находящийся в мертвой стене под индикатором доры, открывается и становится дополнительным индикатором доры, которая называется урадорой. Если в выигранной руке риичи-игрока оказываются урадоры, он получает за них дополнительную стоимость.
	\item \textit{Кан-урадоры}. Если во время раздачи объявлялись каны, при победе после риичи открывается и становится индикатором не только тайл под индикатором доры, но и все тайлы под открытыми индикаторами кандор.
	\item\textit{Акадоры}. В каждой масти одна из пятерок окрашена в красный цвет. такие пятерки также являются дорами, независимо от открытых индикаторов.
\end{itemize}

Если несколько одинаковых индикаторов доры указывают на один и тот же тайл, он становится двойной, тройной или четверной дорой по числу указывающих индикаторов, и такие доры дают соответственно 2, 3 или 4 дополнительных хана за каждый тайл. Если индикатор указывает на пятерку, стоимость акадоры также увеличивается на 1 хан:

\begin{example}

	\mahjong{788m24056p246s22z}

	Мертвая стена:
	\mahjong{XX1z4p1zXX}
	\vspace{5mm}

	Эта рука содержит дор на 7 хан --- по 2 хан за каждый южный ветер, 1 хан за обычную пятерку-дору и 2 хан за пятерку-акадору.

\end{example}

\subsection{Завершение раздачи}

Каждая раздача может завершиться ничьей, абортивной ничьей или победой одного или нескольких игроков.

На последнем тайле раздачи разрешается делать только объявления победы, никакие другие объявления не допускаются.

\subsubsection{Ничья}

\textit{Ничья} (рюкёку) --- это ситуация, когда к концу живой стены ни один игрок не объявил победы.

Все темпай-игроки перед выплатами должны показать другим свои руки, чтобы подтвердить свой темпай. Первым это делает дилер, далее открывают руки все темпай-игроки по ходу игры (юг-запад-север). Если у игрока темпай, но он не хочет его показывать, он также вправе не показывать руку, и тогда ему будет зачтен нотен. Игроки, объявившие риичи, обязаны открыть руку и показать темпай.

При ничьей игроки без темпая (\textit{нотен}-игроки) выплачивают игрокам в темпае суммарно 3000 очков. Если за столом один игрок темпай, три других игрока выплачивают ему по 1000 очков, если два игрока темпай --- каждый из них получает 1500 очков от одного из оставшихся игроков, если три игрока темпая, то единственный нотен-игрок платит каждому по 1000. Если все за столом темпай, или все нотен, никаких выплат не производится.

Ставки риичи при ничьей остаются на столе. Их забирает следующий победитель.

\subsubsection{Победа в раздаче}

Когда один или несколько игроков объявляют победу, раздача немедленно заканчивается и происходит подсчет игровых очков. Игрок, выигравший по цумо, получает выплаты от трех других игроков. Игрок, на сброшенном тайле которого была объявлена победа по рон, выплачивает полную стоимость каждой из выигравших рук.

\subsection{Хонба}

Если в раздаче победил дилер либо была объявлена ничья, справа от дилера кладется \textit{хонба} --- палочка в 100 очков из числа палочек дилера (не вычитается при этом из его счета). Каждая лежащая на столе хонба увеличивает стоимость выигранной в этой раздаче руки на 300 очков.

В случае цумо каждый игрок платит победителю по 100 очков за каждую хонбу. В случае победы по двойному или тройному рону, надбавка за хонбу выплачивается каждому победителю.

\subsection{Смена дилера}

После окончания раздачи определяется, происходит ли смена дилера.
Смена не происходит при выполнении одного из условий:

\begin{itemize}
	\item Дилер является одним из победителей в текущей раздаче
	\item Дилер темпай в случае, если раздача закончилась ничьей.
\end{itemize}

В противном случае происходит смена: Юг становится Востоком, Запад --- Югом, Север --- Западом, Восток --- Севером.

Когда первый дилер снова становится востоком начинается южный раунд.

Когда первый дилер становится востоком в третий раз игра заканчивается.

\subsection{Финальный расчет очков. Ума.}

После завершения южного раунда игра заканчивается. Победителем становится игрок, набравший больше всех очков (также может произойти ничья). Также победитель забирает все оставшиеся на столе палочки риичи (в случае ничьей палочки делятся поровну между победителями).
Затем происходит финальный подсчет:

\begin{itemize}
	\item Из игровых очков каждого игрока вычитаются стартовые очки;
	\item К каждому из результатов добавляется соответствующая месту игрока \textit{ума} (бонус/штраф за занятое место, см. далее);
	\item Если игрок получил штраф во время игры, штраф вычитается из получившегося значения.
\end{itemize}

\subsubsection{Ума}

\textit{Ума} --- это бонусы и штрафы за занятые места.
Бонусы за первое и второе места составляют 15000 и 5000 очков. Штрафы за третье и четвертое места составляют -5000 и -15000 очков соответственно.

Если несколько игроков набрали одинаковое количество очков, то ума за соответствующие места делится между ними поровну.

\subsection{Подсчет стоимости руки и порядок выплат}

Стоимость руки сначала считается как сумма всех хан, полученных за яку и доры, и всех \textit{фу} --- "мини-очков", начисляемых за определенные сеты в руке, форму ожидания на выигрышный тайл и другие особенности руки, не дающие яку. После этого стоимость, выраженная в ханах и фу, переводится в очки по специальной таблице или формуле, которые будут рассмотрены далее. Если в руке 5 хан или больше, фу считать не нужно: стоимость такой руки в очках будет зависеть только от количества хан.

 После подсчета всех фу получившееся число округляется вверх до десятков (например, если при суммировании получилось 54 фу, итоговым значением будет 60). Единственное исключение --- чиитойцу (7 пар): любая рука, содержащая чиитойцу, оценивается в 25 фу.

\subsubsection{Базовые фу}

\begin{tabular}{lc}
	\toprule
	Условие & Фу\\
	\midrule
	Чиитойцу (никакие другие фу не начисляются) & 25 фу\\
	\midrule
	Победа по рон с закрытой рукой & 30 фу\\
	\midrule
	Все остальные случаи & 20 фу\\
	\bottomrule
\end{tabular}

\subsubsection{Фу за поны и каны}

\noindent\begin{tabular}{L{7cm}C{3cm}C{3cm}}
	\toprule
	& Открытый сет &
	Закрытый сет \\
	\midrule
	Пон средних тайлов &
	+ 2 фу &
	+ 4 фу \\
	\midrule
	Пон терминальных или благородных тайлов &
	+ 4 фу &
	+ 8 фу \\
	\midrule
	Кан средних тайлов &
	+ 8 фу &
	+ 16 фу \\
	\midrule
	Кан терминальных или благородных тайлов &
	+ 16 фу &
	+ 32 фу \\
	\bottomrule
\end{tabular}

Если пон образовался в руке с получением выигрышного тайла, он считается открытым, если тайл получен по рону, и закрытым, если по цумо.

\subsubsection{Фу за якухайные пары}

Если в руке пара состоит из тайлов, сет которых дал бы якухай (драконы или ветра раунда/места), за такую пару начисляется 2 дополнительных фу. За пару совпадающих ветров раунда и места, пон которых дал бы двойной якухай, дается 4 фу.

\subsubsection{Фу за ожидание на выигрышный тайл}

\noindent\begin{tabular}{C{3cm}C{5cm}C{3cm}C{3cm}}
\toprule
	Название  &
	Пример формы &
	Выигрышные тайлы &
	Фу \\
\midrule
		\rule[0ex]{0pt}{5ex} Канчан \newline (в "середину") \newline &
	\mahjong{35s} &
	\mahjong{4s} &
	+2 фу \\
\midrule
	\rule[0ex]{0pt}{5ex} Пенчан \newline (в "край") \newline &
	\mahjong{12m} &
	\mahjong{3m} &
	+2 фу \\
\midrule
	\rule[0ex]{0pt}{5ex} Танки \newline (в пару) \newline &
	\mahjong{3s} &
	\mahjong{3s} &
	+2 фу \\
\bottomrule
\end{tabular}

\subsubsection{Дополнительные фу}

\begin{itemize}
	%\item При победе по рону с закрытой рукой дается 10 дополнительных фу;
	\item При победе по цумо с любой рукой (открытой или закрытой) дается 2 дополнительных фу, если эта рука не содержит пинфу\footnote{Можно заметить, что все условия для пинфу подобраны так, чтобы рука в итоге не получила фу кроме базовых 20: в ней нет ни понов и канов, ни якухайных пар, ни ценных ожиданий. Действительно, буквальный перевод названия "пинфу" --- это "рука без фу". Единственное исключение --- +10 фу за рон на закрытой руке.}.
	\item При победе по рону с открытой рукой, в которой выполняются все остальные условия для пинфу (нет понов, канов и якухайных пар, ожидание в чи в две стороны) дается 2 дополнительных фу.
\end{itemize}

\subsection{Подсчет очков в руке}

Для расчетов стоимости руки используют таблицы выплат (см. далее). Используют две таблицы --- одну для расчет рук дилера, вторую для рук не-дилера.

В каждой ячейке сверху большими цифрами указана выплата при победе по рону, ниже --- выплаты каждого из игроков при цумо: в не-дилерской таблице после дроби указана выплата дилера, перед дробью --- каждого из оставшихся игроков.

\subsubsection{Формула перевода ханов и фу в очки}

В таблицах можно заметить закономерность: \textbf{удвоение фу эквивалентно +1 хану}. Эту закономерность полезно знать для лучшего запоминания, кроме того она позволяет быстро рассчитать очки для рук с большим числом фу, например рука стоимостью в 1 хан и 100 фу будет иметь точно ту же стоимость, что и рука в 2 хан 50 фу, и ту же стоимость, что и рука в 3 хан 25 фу. Почему так происходит?

Значения выплат в таблице получены по \textbf{одной и той же формуле}. Во время игры можно пользоваться таблицей, и поэтому знать формулу не нужно, но она может быть полезна для понимания принципов перевода ханов и фу в очки. Выглядит она так:

\begin{equation}
	\large
	Base = Fu \cdot 2 ^ {Han + Bazoro}
\end{equation}

\textit{Бадзоро} --- это базовый коэффициент, обычно равный двум. В данных правилах бадзоро всегда имеет именно такое значение.

Из базового значения легко получить конкретное количество очков по простой формуле:

\begin{equation}
	\large
	Points = \lceil Base * k \rceil
\end{equation}

Обратите внимание на округление вверх до сотен --- оно производится \textit{после} умножения базы на коэффициент \textit{k}.

\textit{k} --- это коэффициент, зависящий от типа победы и дилерства победителя:

\begin{itemize}
	\item При не-дилерском роне: \textit{k=4};
	\item При дилерском роне: \textit{k=6};
	\item При не-дилерском цумо дилер выплачивает победителю сумму, подсчитанную по формуле с \textit{k=2}, а каждый не-дилер --- с \textit{k=1};
	\item При дилерском цумо каждый выплачивает дилеру сумму с \textit{k=2}.
\end{itemize}

Из формулы легко вывести еще несколько закономерностей:

\begin{itemize}
	\item Дилерская рука всегда примерно в 1,5 раза дороже такой же не-дилерской (может быть не совсем точно из-за округления);
	\item Каждый дополнительный хан при неизменном количестве фу увеличивает стоимость руки в 2 раза;
	\item Выплата дилера при не-дилерском цумо равна выплате не-дилера при цумо дилера с такой же рукой;
	\item Суммарная стоимость руки при цумо может быть немного больше стоимости такой же руки при роне из-за того, что при цумо округляется вверх больше чисел. Например, рука в 1 хан 30 фу для не-дилера при роне будет стоить 4*30*8 = 960, округленно 1000 очков, а при цумо будут выплачены суммы в 1*30*8 = 240 (300), еще такие же 300 от другого не-дилера и 2*30*8 = 480 (500) от дилера --- итого получается 1100.
\end{itemize}

\newgeometry{top=1cm,bottom=2cm,left=1cm,right=1cm}
\begin{landscape}
	\small
	\subsubsection{Таблица подсчета очков}
	\noindent\begin{tabular}{|C{2cm}|C{1.7cm}|C{1.7cm}|C{1.7cm}|C{1.7cm}|C{1.7cm}|C{1.7cm}|C{1.7cm}|C{1.7cm}|C{1.7cm}|C{1.7cm}|C{1.7cm}|}
		\hline
		\textbf{Не-дилер} &
		20 фу &
		25 фу &
		30 фу &
		40 фу &
		50 фу &
		60 фу &
		70 фу &
		80 фу &
		90 фу &
		100 фу &
		110 фу \\
		\hline
		1 хан &
		\textbf{-} \linebreak --- &
		\textbf{-} \linebreak --- &
		\textbf{1000} \linebreak 300 / 500 &
		\textbf{1300} \linebreak 400 / 700 &
		\textbf{1600} \linebreak 400 / 800 &
		\textbf{2000} \linebreak 500 / 1000 &
		\textbf{2300} \linebreak 600 / 1200 &
		\textbf{2600} \linebreak 700 / 1300 &
		\textbf{2900} \linebreak 800 / 1500 &
		\textbf{3200} \linebreak 800 / 1600 &
		\textbf{3600} \linebreak 900 / 1800 \\
		\hline
		2 хан &
		\textbf{-} \linebreak 400 / 700 &
		\textbf{1600} \linebreak --- &
		\textbf{2000} \linebreak 500 / 1000 &
		\textbf{2600} \linebreak 700 / 1300 &
		\textbf{3200} \linebreak 800 / 1600 &
		\textbf{3900} \linebreak 1000 / 2000 &
		\textbf{4500} \linebreak 1200 / 2300 &
		\textbf{5200} \linebreak 1300 / 2600 &
		\textbf{5800} \linebreak 1500 / 2900 &
		\textbf{6400} \linebreak 1600 / 3200 &
		\textbf{7100} \linebreak 1800 / 3600 \\
		\hline
		3 хан &
		\textbf{-} \linebreak 700 / 1300 &
		\textbf{3200} \linebreak 800 / 1600 &
		\textbf{3900} \linebreak 1000 / 2000 &
		\textbf{5200} \linebreak 1300 / 2600 &
		\textbf{6400} \linebreak 1600 / 3200 &
		\textbf{7700} \linebreak 2000 / 3900 &
		\multicolumn{5}{c|}{ манган: \textbf{8000}  2000 / 4000 } \\
		\hline
		4 хан &
		\textbf{-} \linebreak 1300 / 2600 &
		\textbf{6400} \linebreak 1600 / 3200 &
		\textbf{7700} \linebreak 2000 / 3900 &
		\multicolumn{8}{c|}{ манган: \textbf{8000}  2000 / 4000 } \\
		\hline
		5 хан &
		\multicolumn{11}{c|}{манган: \textbf{8000}  2000 / 4000} \\
		\hline
		6-7 хан &
		\multicolumn{11}{c|}{ханеман: \textbf{12000} 3000 / 6000} \\
		\hline
		8-10 хан &
		\multicolumn{11}{c|}{ байман: \textbf{16000} 4000 / 8000} \\
		\hline
		11-12 хан &
		\multicolumn{11}{c|}{санбайман: \textbf{24000} 6000 / 12000} \\
		\hline
		13+ хан &
		\multicolumn{11}{c|}{якуман: \textbf{32000} 8000 / 16000} \\
		\hline
	\end{tabular}

	\noindent\begin{tabular}{|C{2cm}|C{1.7cm}|C{1.7cm}|C{1.7cm}|C{1.7cm}|C{1.7cm}|C{1.7cm}|C{1.7cm}|C{1.7cm}|C{1.7cm}|C{1.7cm}|C{1.7cm}|}
		\hline
		\textbf{Дилер} &
		20 фу &
		25 фу &
		30 фу &
		40 фу &
		50 фу &
		60 фу &
		70 фу &
		80 фу &
		90 фу &
		100 фу &
		110 фу \\
		\hline
		1 хан &
		\textbf{-} \linebreak --- &
		\textbf{-} \linebreak --- &
		\textbf{1500} \linebreak 500 &
		\textbf{2000} \linebreak 700 &
		\textbf{2400} \linebreak 800 &
		\textbf{2900} \linebreak 1000 &
		\textbf{3400} \linebreak 1200 &
		\textbf{3900} \linebreak 1300 &
		\textbf{4400} \linebreak 1500 &
		\textbf{4800} \linebreak 1600 &
		\textbf{5300} \linebreak 1800 \\
		\hline
		2 хан &
		\textbf{-} \linebreak 700 &
		\textbf{2400} \linebreak --- &
		\textbf{2900} \linebreak 1000 &
		\textbf{3900} \linebreak 1300 &
		\textbf{4800} \linebreak 1600 &
		\textbf{5800} \linebreak 2000 &
		\textbf{6800} \linebreak 2300 &
		\textbf{7700} \linebreak 2600 &
		\textbf{8700} \linebreak 2900 &
		\textbf{9600} \linebreak 3200 &
		\textbf{10600} \linebreak 3600 \\
		\hline
		3 хан &
		\textbf{-} \linebreak 1300 &
		\textbf{4800} \linebreak 1600 &
		\textbf{5800} \linebreak 2000 &
		\textbf{7700} \linebreak 2600 &
		\textbf{9600} \linebreak 3200 &
		\textbf{11600} \linebreak 3900 &
		\multicolumn{5}{c|}{манган: \textbf{12000} 4000} \\
		\hline
		4 хан &
		\textbf{-} \linebreak 2600 &
		\textbf{9600} \linebreak 3200 &
		\textbf{11600} \linebreak 3900 &
		\multicolumn{8}{c|}{манган: \textbf{12000} 4000} \\
		\hline
		5 хан &
		\multicolumn{11}{c|}{манган: \textbf{12000} 4000} \\
		\hline
		6-7 хан &
		\multicolumn{11}{c|}{ханеман: \textbf{18000} 6000} \\
		\hline
		8-10 хан &
		\multicolumn{11}{c|}{байман: \textbf{24000} 8000} \\
		\hline
		11-12 хан &
		\multicolumn{11}{c|}{санбайман: \textbf{36000} 12000} \\
		\hline
		13+ хан &
		\multicolumn{11}{c|}{якуман: \textbf{48000} 16000} \\
		\hline
	\end{tabular}
\end{landscape}
\restoregeometry

\subsubsection{Лимиты}

Формула выше применяется только для рук, общая стоимость которых по ней получается не выше 8000 очков для не-дилера или 12000 для дилера. Для более дорогих рук стоимости фиксированы и зависят только от числа хан.

\subsubsection{Правило пао}

Правило \textit{Пао} --- это правило, устанавливающее особый порядок выплат за некоторые яку:

\begin{itemize}
	\item Дайсанген. Игрок, на сброшенном тайле которого, оппонент, имеющий два объявленных пона/кана драконов, объявил третий пон/кан драконов, становится "ответственным".
	\item Дайсуши. Игрок, который в процессе раздачи cбросом тайла позволил собрать четвертый пон/кан ветров оппоненту, имеющему три объявленных пона/кана ветров, становится "ответственным".
\end{itemize}

Если после этого якуман был собран, то выплаты производятся следующим образом:

\begin{itemize}
	\item Цумо --- вместо обычных трех выплат "ответственный" игрок платит всю стоимость, как при роне с него;
	\item Рон с другого игрока --- "ответственный" игрок делит с ним выплату пополам.
\end{itemize}

Если условия для пао не выполняются, или же якуман к концу раздачи не собрался (руку собрал другой игрок или произошла ничья), выплаты производятся как обычно.

При выплатах по пао стоимость хонбы выплачивается "ответственным" игроком в случае цумо и сбросившему тайл в рон игроком в случае рона.

Применение правила пао не зависит от того, объявлял ли риичи игрок, сбросивший последний ветер или последнего дракона для якумана.

\subsection{Нарушения и штрафы}

\subsubsection{Мертвая рука}

\textbf{Мертвая рука} --- наказание за нарушения, после которых раздачу возможно продолжить. Игрок, рука которого была объявлена мертвой, до конца раздачи не может делать любые объявления (сетов, победы, риичи). При ничьей мертвая рука считается нотен.

Примеры ситуаций, когда рука игрока объявляется мертвой:

\begin{itemize}
	\item Некорректное объявление сета. Если ошибка обнаружена и исправлена до сноса тайла, мертвая рука не назначается. Пустые объявления "пон", "чи" или "кан", после которых игрок не показывает тайлы, поправляет себя и говорит, что не объявляет сет, не наказываются.
	\item Некорректное объявление победы без открытия руки.
	\item Некорректное объявление риичи (например, не произнесено "риичи", не повернут тайл в дискарде, объявление с открытой рукой).
	\item Замена или попытка замены тайла в руке при риичи.
	\item Нарушение правила \textit{куикаэ}.
	\item Неправильное количество тайлов в руке игрока.
\end{itemize}

\subsubsection{Чомбо}

Чомбо назначается за серьезные нарушения правил, после которых раздача не может продолжаться. Раздача немедленно прерывается и нарушителю назначается штраф в размере двойного шага умы (например, -20000 для умы 15000/5000). Штраф чомбо вычитается из итоговых (после применения умы) очков в конце игры.

После получения кем-либо из игроков чомбо происходит перераздача, при этом дополнительная хонба на стол не выкладывается, риичи-палочки возвращаются владельцам, дилер не меняется.

Примеры ситуаций, когда игроку назначается чомбо:

\begin{itemize}
	\item Некорректное объявление победы с последующим полным открытием руки (объявление победы с рукой без темпая, с рукой без яку, не на выигрышный тайл или рон при действии фуритена / временного фуритена, а также любое объявление кроме "цумо" или "рон").
	\item Нотен-риичи (назначается только в случае ничьей). Если объявление риичи было корректным, но впоследствии рука была объявлена мертвой чомбо не назначается, а игрок в случае ничьей платит обычный штраф за нотен.
	\item Некорректное объявление закрытого кана после риичи (назначается только если игрок объявляет победу, либо в случае ничьей)
	\item Закрытие руки при ничьей после объявления риичи.
	\item Попытка любого объявления в случае если рука была объявлена мертвой.
	\item В случае если игра не может продолжаться из-за ошибки игрока.
\end{itemize}

Чомбо может применяться к нескольким игрокам одновременно.

В случае, если чомбо происходит одновременно с победой другого игрока, засчитывается победа, а штраф чомбо не назначается.

\section{Cписок яку}

Для некоторых яку необходимым условием является закрытая рука. При выигрыше по рон рука остается закрытой.

Все яку могут сочетаться друг с другом, если не сказано обратное. Якуманы сочетаются только между собой.

\subsection{1 хан}

\textbf{Риичи} --- яку за соответствующее объявление.

\textbf{Дабл-риичи} --- дополнительное яку за объявление "риичи" на первом ходу игрока. На момент объявления последовательность ходов не должна быть прервана объявлениями сетов.

\textbf{Иппацу} --- дополнительное яку за победу на первом круге после объявления "риичи". Круг не должен быть прерван объявлениями сетов.

\textbf{Мензен цумо} --- выигрыш со стены с закрытой рукой.

\textbf{Пинфу} --- закрытая рука, состоящая только из чи и пары, не дающей фу. Выигрышный тайл должен завершать чи с ожиданием в две стороны.

\textbf{Иипейко} --- закрытая рука с двумя одинаковыми (одинакового достоинства и масти) чи. Не сочетается с яку "рянпейко" и "чиитойцу".

\mahjong{223344p}

\textbf{Таняо} --- рука без терминальных и благородных тайлов.

\textbf{Саншоку додзюн} или просто \textbf{саншоку} --- рука с тремя чи одинакового достоинства в разных мастях. +1 хан в закрытой руке.

\mahjong{123p 123s 123m}

\textbf{Иццу} --- рука с тремя последовательными чи одной масти, составляющие ряд от 1 до 9. +1 хан в закрытой руке.

\mahjong{123s 456s 789s}

\textbf{Якухай} --- пон/кан драконов, ветров места или ветров раунда. Суммируются между собой.

\textbf{Чанта} --- рука с хотя бы одним терминальным или благородным тайлом в каждом из сетов и парой терминальных или благородных тайлов. Рука должна содержать благородные тайлы и минимум 1 чи. +1 хан в закрытой руке. Не сочетается с яку "джунчан".

\textbf{Риншан кайхо} --- победа на тайле замены после объявления кана. Считается победой со стены ("цумо").

\textbf{Чанкан} --- победа на тайле, которым другой игрок дополняет свой открытый пон до кана. Кан-дора в этом случае не открывается. Может сочетаться с иппацу, если перед этим не было других объявлений сетов.

\textbf{Хайтей раюоэ} или просто \textbf{хайтей} --- победа на последнем тайле, взятом из стены. Не сочетается с "Риншан кайхо".

\textbf{Хотей раоюй} или просто \textbf{хотей} --- победа на последнем сброшенном тайле в раздаче.

\subsection{2 хан}

\textbf{Чиитойцу} --- рука, в которой содержится семь различных пар. Всегда стоит \textbf{25} фу. Не сочетается с яку "иипейко" и "рянпейко".

\textbf{Саншоку доко} --- рука с тремя понами/канами одинакового достоинства в разных мастях.

\mahjong{444m 444s 444p }

\textbf{Сананко} --- рука с тремя закрытыми понами/канами. Рука при этом может быть открытой.

\textbf{Санканцу} --- рука с тремя канами.

\textbf{Тойтой} --- рука с четырьмя понами/канами.

\textbf{Хоницу} --- рука из тайлов одной масти и благородных. +1 хан в закрытой руке.

\textbf{Шосанген} --- рука с двумя понами/канами драконов и парой драконов. Сочетается с "якухай".

\textbf{Хонрото} --- рука из терминальных и благородных тайлов.

\textbf{Джунчан} --- рука с хотя бы одним терминальным тайлом в каждом из сетов и парой терминальных тайлов. Рука должна содержать минимум одно чи. +1 хан в закрытой руке.

\subsection{3 хан}

\textbf{Рянпейко} --- закрытая рука с двумя парами одинаковых чи. Не сочетается с яку "иипейко" и "чиитойцу" .

\mahjong{112233p 556677s 5m}

\subsection{5 хан}

\textbf{Чиницу} --- рука из тайлов одной масти. +1 хан в закрытой руке.

\textbf{Ренхо} --- победа со сброса на первом круге до того как игрок сделал свой первый ход. Круг не должен быть прерван объявлениями сетов.

\subsection{Якуманы}

\textbf{Кокуши мусо} или \textbf{13 сирот} --- рука, в которой содержатся все тринадцать терминальных и благородных тайлов плюс пара к любому из них. Игрок может объявить рон при объявлении противником закрытого кана выигрышных тайлов. Это единственный случай, когда возможен рон тайлом из закрытого кана.

\textbf{Чууренпото} ---закрытая рука, содержащая тайлы 1112345678999 одной масти и еще один любой тайл той же масти.

\mahjong{1112345678999m} \hfill \mahjong{7m}

\textbf{Тенхо} --- победа дилера сразу с набранной со стены рукой в начале раздачи.

\textbf{Чихо} --- победа не-дилера со стены на первом круге. Круг не должен быть прерван объявлениями сетов.

\textbf{Сууанко} --- четыре закрытых пона/кана.

\textbf{Сууканцу} --- четыре кана.

\textbf{Рюисо} --- рука, содержащая только зеленые тайлы --- 23468 в масти соу и зеленого дракона. Наличие зеленых драконов необязательно.

 \textbf{Чинрото} --- рука из терминальных тайлов.

 \textbf{Цуисо} --- рука из благородных тайлов.

 \textbf{Дайсанген} --- три пона драконов.

\textbf{Шосуши} --- три пона ветров и пара ветров.

\textbf{Дайсуши} --- четыре пона ветров в руке.

\section{Вариации правил}

Данный раздел необязателен к прочтению, однако может быть полезен в случаях, когда на турнире заявляются некоторые отличия от набора правил по умолчанию. Все вариации, описанные ниже, не являются вариациями по умолчанию, но разрешены при явном указании в регламенте турнира.

\subsection{Без акадор}

Существует вариация правил, в которой акадоры отсутствуют. Красные пятерки в наборах будут заменены обычными. Дополнительных правил в этом случае не вводится.

В правилах по умолчанию: акадоры есть.

\subsection{Старт с 25000}

В клубных и онлайн-играх может быть принято начинать игру с 25000 очков. Часто в этом случае также устанавливается ока (бонус первому месту) в размере 20000 очков. На турнирах допускается начинать игру с 25000 очков, однако введение бонуса ока не допускается.

В правилах по умолчанию: старт с 30000.

\subsection{Ума}

Ума --- ранговый бонус или штраф в зависимости от занятого места. Когда говорят, что ума составляет 15000 / 5000, подразумевают, что:

\begin{itemize}
    \item 1 место получит +15000
    \item 2 место получит +5000
    \item 3 место получит -5000
    \item 4 место получит -15000
\end{itemize}

Предпочтительная ума на турнирах - 15000 / 5000. При этом допускается любая другая \textit{симметричная} ума (например, 30000 / 10000). Под симметричной понимается такая ума, которая дает равномерный шаг между местами, т.е. например в случае умы 30000 / 10000 разница бонуса между первым и вторым, вторым и третьим, третьим и четвертым местом одинакова и составляет 20000 очков.

В правилах по умолчанию: ума 15000 / 5000.

\subsection{Нагаши манган}

\textbf{Нагаши манган} --- особая игровая ситуация при ничьей и выполнении следующих условий:

\begin{itemize}
	\item Живая стена полностью исчерпана;
	\item У кого-то из игроков в дискарде лежат исключительно терминальные и благородные тайлы;
	\item Из дискарда этого игрока не было взято в сет ни одного тайла;
	\item Этому игроку не была назначена мертвая рука.
\end{itemize}

Выигравшему игроку (или игрокам, если их несколько) выплачивается стоимость мангана по цумо (т.е. 4000 с дилера и по 2000 с остальных, либо по 4000 с каждого в случае если нагаши манган собрал дилер). Стандартные выплаты по ничьей при этом не производятся.

Нагаши манган не считается выигрышной рукой. Это особая игровая ситуация (разновидность ничьей с иной механикой выплат), поэтому никакие другие механики, относящиеся к победам, не применяются, в частности:

\begin{itemize}
	\item Риичи-палочки на столе остаются на месте (риичи, объявленные в ходе раздачи, остаются на столе и разыгрываются в следующей раздаче);
	\item Дополнительных выплат за хонбы не производится;
	\item Смена дилера и добавление хонбы происходит исходя из того, остался ли дилер темпай в конце раздачи (аналогично обычной ничьей).
\end{itemize}

Игрок, собравший нагаши манган, не обязан иметь темпай к концу раздачи.

В случае, если по итогу раздачи назначается нагаши манган, но кто-либо из игроков должен получить штраф (например, за нотен риичи), приоритет имеет штраф.

В правилах по умолчанию: нагаши манган отсутствует.

\subsection{Банкротство}

\textit{Тоби} (\textit{банкротство}) --- правило, согласно которому игра заканчивается досрочно, как только один или несколько игроков уходят в минус по очкам. Кроме того, оно не позволяет делать ставку риичи, если у игрока в наличии менее 1000 очков.

В правилах по умолчанию: банкротство отсутствует.

\subsection{Абортивные ничьи}

\textbf{Абортивная ничья} --- это ничья, которая происходит при выполнении одного из следующих условий:

\begin{itemize}
	\item Все игроки сбросили один и тот же ветер на первом круге раздачи. Круг не должен быть прерван объявлениями сетов.
	\item Ситуация \textbf{кюсюкюхай} --- на первом ходе игрока у него в руке есть 9 или более уникальных терминальных и благородных тайлов. В этом случае игрок может открыть эти 9 уникальных тайлов из руки и объявить абортивную ничью. До объявления кюсюкюхай круг также не должен быть прерван объявлениями сетов, в ином случае возможность объявить абортивную ничью теряется.
	\item Все игроки за столом объявили риичи, и на последнем объявлении риичи никто не объявил победу.
	\item Было открыто 4 кана разными игроками и на последнем сбросе после объявления никто не объявил победу. В случае если все четыре кана объявлены одним и тем же игроком, игра продолжается.
\end{itemize}

Существует также вариация правил, в которой при объявлении победы тремя игроками с одного и того же сброшенного тайла также объявляется абортивная ничья.

В случае любой абортивной ничьей все игроки, объявившие риичи, обязаны показать свой темпай. Если у кого-то обнаруживается нотен-риичи, назначается штраф чомбо.

В правилах по умолчанию: абортивные ничьи отсутствуют.

\subsection{Атамаханэ}

Правило атамаханэ говорит о том, что при победе по рон с одного и того же сброшенного тайла приоритет имеет победа игрока, сидящего первым по ходу игры от сбросившего победный тайл. Победа следующих по кругу игроков не засчитывается. Все ставки риичи со стола забирает игрок, победа которого была засчитана.

В правилах по умолчанию: атамаханэ отсутствует.

\subsection{Чомбо как обратный манган}

В клубных играх может использоваться вариант чомбо как "обратного мангана" --- это выглядит как "цумо наоборот": если нарушитель --- дилер, он платит по 4000 очков каждому, а если нет, то 4000 очков дилеру и по 2000 двум остальным игрокам, в этом случае дополнительный штраф после умы не применяется.

В правилах по умолчанию: чомбо как обратный манган отсутствует.

\subsection{Якитори}

Правило якитори добавляет на игровой стол дополнительный элемент --- пластинку с изображением птички возле каждого игрока. Если игрок побеждает, он переворачивает либо убирает со стола пластинку. В конце игры те игроки, у которых пластинка осталась неперевернутой или осталась на столе, получают штраф в размере двойного шага умы (20000 в случае умы 15000 / 5000 или 40000 в случае умы 30000 / 10000). Это правило стимулирует игроков стремиться победить любой ценой хотя бы раз за игру, однако вносит существенные нюансы в игровую механику и потому по умолчанию не рекомендуется.

В правилах по умолчанию: якитори отсутствует.

\subsection{Отсутствие ренхо}

Допускается исключить яку ренхо с турнира. В этом случае игрок не имеет права победить с другого игрока при отсутствии других яку на руке на первом непрерванном круге игры.

В правилах по умолчанию: ренхо присутствует и стоит манган.
